\documentclass[11pt]{article}

\usepackage[T1]{fontenc}
\usepackage[utf8]{inputenc}
\usepackage[english]{babel}
\usepackage{amsmath,amssymb,amsthm}
\usepackage{booktabs}
\usepackage{geometry}
\usepackage{hyperref}
\usepackage{url}
\geometry{margin=1in}

\hypersetup{
  colorlinks=true,
  linkcolor=black,
  citecolor=blue,
  urlcolor=blue,
  pdftitle={Exact Extension of OEIS A049230 through 27 Steps},
  pdfauthor={Carlo Corti}
}

\newtheorem{theorem}{Theorem}
\newtheorem{proposition}[theorem]{Proposition}
\newtheorem{computationalresult}[theorem]{Computational Result}
\theoremstyle{definition}
\newtheorem{definition}[theorem]{Definition}
\theoremstyle{remark}
\newtheorem{remark}[theorem]{Remark}

\title{Exact Extension of OEIS A049230 through 27 Steps:\\
Cubic-Orbit Enumeration at 19--22 and Published Partition-Function Data at 23--27}
\author{Carlo Corti\\
Independent researcher\\
\texttt{carlo.corti@outlook.com}}
\date{August 2026}

\begin{document}
\maketitle

\begin{abstract}
OEIS A049230 counts rooted $n$-step self-avoiding walks on the simple cubic
lattice having exactly two nearest-neighbour contacts between vertices that
are nonconsecutive along the walk.
Translations are removed by fixing the initial vertex, whereas rotations,
reflections, and reversal remain distinct.  At the time of our audit, OEIS
A049230 listed a contiguous prefix through $a(18)$.  Our literature search
identified exact published partition-function coefficients that normalize to
$a(23),\ldots,a(27)$.  Relative to the OEIS prefix, this left precisely four
intervening values undetermined.  We compute them as
\[
\begin{aligned}
a(19)&=2126459849880,&
a(20)&=9600694064496,\\
a(21)&=43090682243160,&
a(22)&=192553331372616.
\end{aligned}
\]
The computation enumerates genuinely three-dimensional walks modulo the full
48-element cubic point group and reconstructs the rooted count by the exact
identity
\[
 a(n)=3\,A033323(n)+48p_{n,2}^{(3)}.
\]
On the computational platform specified below, the frozen campaign used 2411
deterministic recovery units in three segments and records 10232.08 seconds of
accumulated per-unit wall-clock time.  An
independent Java implementation directly reproduced $a(19)$; it differs in
software realization and symmetry handling but shares the same exhaustive
depth-first, incremental-contact method.  The normalized published values at
$n=23,\ldots,27$ are not outputs of the campaign.  The resulting documented
range is contiguous for $1\leq n\leq27$; within the evidence assembled in the
audit of 26 August 2026, the first unresolved index is $n=28$.
\end{abstract}

\noindent\textbf{2020 Mathematics Subject Classification.}
Primary 05A15; Secondary 82B41, 68W10.

\noindent\textbf{Key words and phrases.}
self-avoiding walk, interacting self-avoiding walk, exact enumeration,
nearest-neighbour contact, cubic symmetry, OEIS.

\section{Introduction and provenance}

A self-avoiding walk (SAW) is a nearest-neighbour lattice path that visits no
vertex twice.  Contact-resolved enumeration is a classical finite-lattice model
of polymer--solvent interaction~\cite{FisherHiley1961}.  Nemirovsky, Freed,
Ishinabe, and Douglas organized these counts by the number of coordinate
dimensions actually spanned and introduced dimension-independent coefficients
$p_{n,m}^{(l)}$~\cite{NemirovskyEtAl1992}.  Their framework fixes the symmetry
normalization used below.

OEIS A049230 was submitted by N.~J.~A. Sloane.  Petros Hadjicostas clarified
its contact interpretation on 3 January 2019, and Sean A. Irvine supplied the
displayed terms $a(12),\ldots,a(18)$ on 23 July 2021~\cite{OEISA049230}.  A
live inspection on 26 August 2026 still displayed the contiguous range only
through $a(18)$.  During our literature audit, we identified exact
partition-function papers containing later isolated coefficients corresponding
to $a(23),\ldots,a(27)$~
\cite{LeeEtAl2012,HsiehEtAl2016EPJ,HsiehEtAl2016CPC}; the later study of
Huang, Hsieh, and Chen independently confirms the exact-enumeration range
through 27 steps, while referring to those partition-function sources for
fixed-first-step partition polynomials~\cite{HuangEtAl2022}.  After exact
normalization, this made $n=19,20,21,22$ the only intervening unresolved
indices, rather than every index after 18.

This paper reports an exhaustive computation of precisely those four values.
It also normalizes the later published coefficients, thereby documenting a
contiguous exact range through 27 without presenting published values as new
campaign outputs.  The distinction matters both scientifically and for the
provenance of the associated b-file.

\section{Definition, indexing, and dimensional decomposition}

\begin{definition}
Let $w=(v_0,v_1,\ldots,v_n)$ be an ordered nearest-neighbour walk on
$\mathbb Z^3$, rooted by $v_0=0$.  It is self-avoiding if its vertices are
distinct.  Its number of contacts is
\[
 m(w)=\#\bigl\{\{i,j\}:0\leq i<j\leq n,\ j-i>1,
                   \ \lVert v_i-v_j\rVert_1=1\bigr\}.
\]
The value $a(n)=A049230(n)$ is the number of such walks with $m(w)=2$.
The vertex order is retained, so reversal is distinct; no rotation or
reflection quotient is taken in $a(n)$.
\end{definition}

The walk has $n$ steps and $N=n+1$ monomers.  This conversion is essential
when partition polynomials are indexed by monomer count.  In this paper,
\[
 Z_N(x)=\sum_{w\in\mathcal W_{N-1}^{(3)}}x^{m(w)},
\]
where $\mathcal W_{N-1}^{(3)}$ is the set of rooted ordered $(N-1)$-step
simple-cubic SAWs with all six possible first-step directions.  Let
$\widehat{\mathcal W}_{N-1}^{(3)}$ be the subset whose first step is $+x$, and
write
\[
 \widehat Z_N(x)=
 \sum_{w\in\widehat{\mathcal W}_{N-1}^{(3)}}x^{m(w)}.
\]
The cubic point group acts transitively on the six first-step directions
without changing contacts, so
\[
 Z_N(x)=6\widehat Z_N(x),\qquad
 a(n)=[x^2]Z_{n+1}(x)=6[x^2]\widehat Z_{n+1}(x).
\]

\begin{definition}
For $d\geq1$, let $C_{n,m}(d)$ be the number of rooted ordered $n$-step SAWs
on $\mathbb Z^d$ with exactly $m$ contacts, with all $2d$ first-step
directions retained and with reversal distinct.  For $1\leq l\leq n$, let
$p_{n,m}^{(l)}$ be the number of orbits of rooted ordered $n$-step SAWs on
$\mathbb Z^l$ with exactly $m$ contacts whose step vectors use all $l$
coordinate axes, under the order-$2^l l!$ hyperoctahedral group of signed
coordinate permutations of $\mathbb Z^l$.  These walks are called genuinely
$l$-dimensional; lower-dimensional walks are excluded from $p_{n,m}^{(l)}$.
Embedding $\mathbb Z^l$ in any fixed coordinate $l$-subspace of $\mathbb Z^d$
gives a bijection with the genuinely $l$-dimensional walks in that subspace.
Thus the orbit count is independent of the ambient dimension, while the
choice of subspace contributes the factor $\binom dl$ below.
\end{definition}

Nemirovsky et al.\ give the dimensional decomposition
\begin{equation}
 C_{n,m}(d)=\sum_{l=1}^{\min(n,d)}2^l l!\binom dl p_{n,m}^{(l)}.
 \label{eq:dimension}
\end{equation}
This is Eq.~(5) on p.~1090 of their paper~\cite{NemirovskyEtAl1992}.

\begin{proposition}
For every $n\geq1$,
\begin{equation}
 a(n)=3\,A033323(n)+48p_{n,2}^{(3)}.
 \label{eq:reconstruction}
\end{equation}
\end{proposition}

\begin{proof}
A one-dimensional SAW is a monotone straight walk, so it has no contact
between nonconsecutive vertices and $p_{n,2}^{(1)}=0$.  OEIS A033323 uses the
same step index, rooting, vertex order, and contact convention on the square
lattice; hence
\[
 A033323(n)=C_{n,2}(2)=2^2 2!\,p_{n,2}^{(2)}=8p_{n,2}^{(2)}.
\]
Specializing Eq.~\eqref{eq:dimension} to $d=3$ gives
\[
 C_{n,2}(3)=24p_{n,2}^{(2)}+48p_{n,2}^{(3)}
           =3A033323(n)+48p_{n,2}^{(3)}.
\]
Since $C_{n,2}(3)=a(n)$, the result follows.
\end{proof}

The derivation is structural, not a numerical fit, and the two input tables
were joined on the mathematical index $n$.  In particular, the planar inputs
for $n=19,\ldots,22$ originate in Table B1, pp.~4738--4739, of Bennett-Wood
et al.\ (Table 4, pp.~26--28, in the corresponding preprint).  Its column
$c_n(2)/4$ is $A033323(n)/4$ in the present notation, so multiplication by
four recovers $A033323(n)$~\cite{BennettWoodEtAl1998,OEISA033323}.  For
$n=19,\ldots,22$, the four recovered values coincide with the corresponding
OEIS A033323 terms.  The OEIS record
credits Hadjicostas with the source interpretation, Irvine with terms
$a(20),\ldots,a(26)$, and Corti with $a(27),\ldots,a(29)$; the primary
numerical provenance of all planar values used here remains the
Bennett-Wood et al.\ table.

\section{Symmetry and stabilizers}

The full octahedral group $O_h$, equivalently the three-dimensional
hyperoctahedral group, consists of the $2^3 3!=48$ signed coordinate
permutations.  It acts on a rooted walk coordinatewise.  Its action is free on
a genuinely three-dimensional ordered walk: an element fixing the ordered
walk fixes every vertex, hence every step vector; because the walk contains
three linearly independent step vectors, the element is the identity.
Consequently every genuinely three-dimensional orbit has size 48.

Planar walks behave differently.  Reflection through their supporting
coordinate plane fixes every vertex, so they have a nontrivial stabilizer in
the three-dimensional group.  Their contribution is therefore not obtained by
blindly dividing the full count by 48.  Instead A033323 supplies the oriented
count in one abstract square lattice and the factor 3 in
Eq.~\eqref{eq:reconstruction} chooses one of the three coordinate planes.  This
separation is why the reconstruction treats planar and genuinely
three-dimensional strata independently.

\begin{proposition}
Every genuinely three-dimensional orbit has exactly one representative that
satisfies all of the following conditions:
\begin{enumerate}
\item the first step is $+x$;
\item the first step along the $y$-axis is $+y$, and no step along the $z$-axis
      occurs before it;
\item the first step along the $z$-axis is $+z$.
\end{enumerate}
\label{prop:canonical}
\end{proposition}

\begin{proof}
The first condition chooses one of the six signed directions of the first
step.  After it is fixed, the residual group has order eight.  Let $t_y$ and
$t_z$ be the first times at which the image uses the $y$- and $z$-axes;
genuine three-dimensionality makes both finite and distinct.  Of the identity
and the transverse-axis transposition, exactly one gives $t_y<t_z$.  The two
remaining sign freedoms then independently make the first $y$ step
$+y$ and the first $z$ step $+z$.  Hence the transverse-axis ordering removes
the residual swap, the two sign conditions remove the residual reflections,
and all eight elements of the stabilizer of $+x$ are resolved uniquely.
\end{proof}

The frozen production search implements precisely these conditions: it fixes
the first step to $+x$, rejects a $z$ step until a positive first $y$ step has
occurred, rejects a negative first step along either transverse axis, and
counts a state only after $z$ has been used.  Proposition~
\ref{prop:canonical} therefore implements the factor 48 in
Eq.~\eqref{eq:reconstruction}; reversal is never identified.

The corresponding partial-walk admissibility predicate is closed under taking
prefixes: a prefix of a canonical walk cannot have the wrong first step, use
$z$ before $y$, or have a negative first use of either transverse axis; an axis
not yet used creates no violation.  Since the contact count is also
nondecreasing, every canonical target walk of length at least 10 has an
admissible length-10 prefix with at most two contacts.  The predicate may
therefore be enforced incrementally, and the depth-10 frontier is exhaustive.

\section{Exhaustive enumeration}

The enumerator is a single C17/OpenMP program.  For a maximum length $n$, each
worker maintains a private byte-valued occupancy grid of side $2n+3$, padded
by one lattice layer beyond every reachable coordinate.  It fixes the first
step to $+x$ as in Proposition~\ref{prop:canonical}, then extends the walk in
the deterministic direction order $+x,-x,+y,-y,+z,-z$.  A candidate vertex is
rejected if it is already occupied.  Otherwise the number of newly created
nonbonded contacts is
\[
 \#\{\hbox{occupied nearest neighbours of the candidate}\}-1,
\]
because the predecessor is the unique occupied bonded neighbour.  Contacts
cannot disappear as a branch grows, so any branch exceeding two contacts is
pruned.  At every depth $n$, a genuinely three-dimensional state with exactly
two contacts contributes one to $p_{n,2}^{(3)}$.

The implementation, build recipe, campaign outputs, and validation materials
described in this paper form the accompanying project package~
\cite{CortiPackage2026}.

The search tree was split at depth 10 into 154271 deterministic prefix tasks,
where a task is one admissible canonical partial walk together with its
contact and transverse-axis state.  The tasks were grouped in generation order
into 2411 recovery units, each a consecutive range of at most 64 tasks.  A
segment is one operator invocation that completes whole recovery units until
the configured time target or the final unit is reached.  The manifest is the
durable table containing the frozen configuration, segment state, and one
ordered record per completed unit, including its task interval, elapsed time,
and per-depth counts.

Each OpenMP worker used a private occupancy grid and private
\texttt{uint64\_t} counters.  A critical merge performed checked additions and
failed on overflow; the program's structural cap at $n=29$ is justified by the
nonbacktracking orbit bound $p_{n,2}^{(3)}\leq5^{n-1}/8$, which is below the
64-bit limit through that cap.  The manifest bound the frozen source and build
identity to the configuration, recorded every completed unit exactly once,
and published state by no-clobber atomic replacement.  The verifier regenerated
the depth-10 frontier and required its task count to equal the manifest's
154271 before aggregating all unit records.  Interrupted execution resumes
only at a recovery-unit boundary, so segmentation changes neither the search
space nor the arithmetic result.

The official configuration used 8 OpenMP threads, split depth 10, unit size 64,
and a 4500-second soft segment target.  The target computer had an AMD Ryzen 9
7940HS processor (8 cores, 16 logical processors) and 64 GB DDR5 installed RAM;
the operating system reported 61 GiB total memory.  It ran Linux Mint 22.3
with 64-bit Linux kernel 6.17.0-42-generic.  The executable was built with GCC
13.3.0 and GNU libgomp.  The relevant Makefile flags were
\begin{quote}
\small\ttfamily
-O3 -flto -funroll-loops -fomit-frame-pointer -march=native\par
-mtune=native -std=c17 -Wall -Wextra -Wpedantic -Wshadow\par
-Wconversion -Wdouble-promotion -Wformat=2 -fopenmp -lm
\end{quote}

The campaign was run against that frozen software and environment.
It completed all 2411 units in three normally terminating segments.
Their individually rounded compute times were 4500, 4501, and 1231 seconds.  These displayed values sum to 10232 seconds.  The target is checked
only before starting the next complete recovery unit, which explains the
small overshoot of the first two segments.  These quantities are wall-clock
compute time, not CPU time or campaign calendar time; manual pauses between
segments are excluded.  They are specific to the stated hardware, software,
compiler flags, and eight-thread configuration.

\section{Validation and independence boundary}

Validation addressed the mathematical predicate, symmetry normalization,
indexing, persistence, and implementation separately.

\begin{enumerate}
\item The frozen executable reproduced $p_{n,2}^{(3)}$ for $1\leq n\leq18$.
  Reconstruction by Eq.~\eqref{eq:reconstruction} reproduced the previous
  A049230 prefix byte for byte.
\item A direct oriented oracle exhaustively enumerated the full rooted target
  set through $n=8$ without using the production canonical-orientation
  predicate.  A separate signed-permutation oracle partitioned the same full
  target sets into orbits, explicitly generated all 48 images of every orbit,
  and verified both orbit size 48 and exactly one representative satisfying
  Proposition~\ref{prop:canonical}.
\item Crippen's independent Table I gives the number $N_2(n+1)$ of cubic
  symmetry classes with two contacts for $5\leq n\leq11$.  In the present
  notation it checks
  \[
    N_2(n+1)=p_{n,2}^{(2)}+p_{n,2}^{(3)}
            =\frac{a(n)+3A033323(n)}{48},
  \]
  after reconciling rooting, reversal, and symmetry conventions~
  \cite{Crippen2000}.  For $n=5,\ldots,11$, the common values are
  $8,52,270,1446,7578,38473,191154$.
\item The primary Shirai--Sakumichi attachment tabulates endpoint-constrained
  counts $W_{n,m}(r)$ for endpoints $(r,0,0)$.  Strict aggregation at $m=2$
  gave $\sum_rW_{19,2}(r)=6927757800$ and
  $\sum_rW_{20,2}(r)=28720091852$; multiplying by the six axial directions
  gives exact subsets 41566546800 and 172320551112~
  \cite{ShiraiSakumichi2023}.  These marginals test depth-19 and depth-20
  coverage but are not the unrestricted totals.
\item The completed manifest passed strict continuity, integrity, source
  identity, task-coverage, and final-state checks.  Thread-count and split-mode
  equivalence tests, compiler warnings, two compiler families, static
  analyzers, address/undefined-behaviour sanitizers, and Valgrind checks were
  recorded before the campaign.  Unavailable ThreadSanitizer and system-call
  tracing environments were recorded as unavailable, not as passes.
\end{enumerate}

After the result was frozen, a separate Java verifier based on Sean A.
Irvine's jOEIS class \texttt{ExactContactsWalker} fixed the first step to
$+x$, directly enumerated that complete rooted sector, and restored the six
first-step directions by exact weighting.  It returned 2126459849880 at
$n=19$.  It does not use the production cubic-orbit selector, so this run also
checks the factor-48 normalization at that index.  The verifier was bound to a
fixed jOEIS source revision recorded in the validation package.  The upstream
sources and local harness were verified, compiled in a fresh temporary
directory, and did not read the C campaign, A033323, the Python derivator, or
the b-file~\cite{JoeisRepository}.

This is an \emph{independent software reproduction}, not an independent
mathematical method.  The two implementations differ in author, language,
coordinate and occupancy representation, symmetry accounting, parallel
decomposition, compiler chain, and result path.  Both nevertheless perform
exhaustive depth-first SAW enumeration with incremental contact accumulation.
The validation claim is deliberately limited to $a(19)$.

\section{Results and provenance boundary}

The computed orbit counts and exact reconstructions are shown in
Table~\ref{tab:computed}.  The planar term is small but was retained exactly;
no rounded or row-position arithmetic entered the derivation.

\begin{table}[ht]
\centering
\caption{New campaign results.  The final column is evaluated exactly by
Eq.~\eqref{eq:reconstruction}.}
\label{tab:computed}
\begin{tabular}{r@{\quad}r@{\quad}r@{\quad}r}
\toprule
$n$ & $p_{n,2}^{(3)}$ & $3A033323(n)$ & $a(n)$ \\
\midrule
19 & 44296694918   & 218493816  & 2126459849880 \\
20 & 200002700951  & 564418848  & 9600694064496 \\
21 & 897692330931  & 1450358472 & 43090682243160 \\
22 & 4011450432806 & 3710597928 & 192553331372616 \\
\bottomrule
\end{tabular}
\end{table}

\begin{computationalresult}
Subject to correct execution of the frozen C17/OpenMP enumerator and the
integrity checks on its completed manifest, the exhaustive campaign determines
the four A049230 values in Table~\ref{tab:computed}.  The independently
implemented jOEIS run reproduces the first of these values.
\end{computationalresult}

The later rows have a different provenance.  Lee, Kim, and Lee print the full
$N=24$ contact inventory with all first-step directions in Table I on p.~2;
its contact-number row $K=2$ is 856234452257592, hence $a(23)$ directly~
\cite{LeeEtAl2012}.  Hsieh, Chen, and Hu print fixed-first-step polynomials,
denoted here by $\widehat Z_N$: $\widehat Z_{27}$ and $\widehat Z_{26}$ on p.~4
and $\widehat Z_{25}$ on p.~5 of their EPJ paper, and $\widehat Z_{28}$ in
Sec.~3.1 on p.~6 of their CPC paper.  Their $x^2$ coefficients for
$\widehat Z_{25}$, $\widehat Z_{26}$, $\widehat Z_{27}$, and
$\widehat Z_{28}$ are, respectively,
\[
632297080298076,
2790018300205392,
12274198909485704,
53807607703672080.
\]
Multiplication by six and conversion $n=N-1$ give $a(24),\ldots,a(27)$~
\cite{HsiehEtAl2016EPJ,HsiehEtAl2016CPC}.  Huang, Hsieh, and Chen later use
exact end-to-end polynomials through 27 steps and explicitly direct readers to
these earlier sources for the corresponding fixed-first-step polynomials~
\cite[Appendix A]{HuangEtAl2022}.  Table~\ref{tab:chain} keeps the published
data visibly separate from the campaign.

\begin{table}[ht]
\centering
\caption{The completed exact chain.  ``Campaign'' means exhaustive output of
the computation reported here; ``published'' means normalized primary-source
data, not a campaign output.}
\label{tab:chain}
\begin{tabular}{r@{\quad}r@{\quad}l}
\toprule
$n$ & $a(n)$ & Evidence class \\
\midrule
19 & 2126459849880       & campaign; Java reproduction \\
20 & 9600694064496       & campaign \\
21 & 43090682243160      & campaign \\
22 & 192553331372616     & campaign \\
23 & 856234452257592     & published, Lee--Kim--Lee \\
24 & 3793782481788456    & published, $6[x^2]\widehat Z_{25}$ \\
25 & 16740109801232352   & published, $6[x^2]\widehat Z_{26}$ \\
26 & 73645193456914224   & published, $6[x^2]\widehat Z_{27}$ \\
27 & 322845646222032480  & published, $6[x^2]\widehat Z_{28}$ \\
\bottomrule
\end{tabular}
\end{table}

The canonical b-file contains exactly the continuous index range $1$ through
27.  Rows 1--18 are the unchanged prior prefix, rows 19--22 are the campaign
reconstructions, and rows 23--27 are the normalized published coefficients.
A dated literature and exact-data audit completed on 26 August 2026 searched
the A-number, the notations $C_{n,2}(3)$ and $p_{n,2}^{(3)}$, cubic contact
classes, exact $Z_N(x)$ coefficients, related OEIS records, forward and
backward citations of the primary papers, supplementary data, and exact
integer fingerprints.  An admissible datum was required to be a full
unrestricted rooted two-contact count with enough convention information for
exact normalization.  No such coefficient was found for $n=28$ or above.
Accordingly, within the evidence assembled in this audit, 27 is the documented
endpoint and 28 the first unresolved index.  This bounded statement neither
claims universal nonexistence nor authorizes a recomputation of a later row.

\section{Finite numerical observations}

Two exact finite patterns provide useful checks but are not asymptotic claims.
First, Eq.~\eqref{eq:reconstruction} implies divisibility by 24: both
$3A033323(n)=24p_{n,2}^{(2)}$ and $48p_{n,2}^{(3)}$ are multiples of 24.
Every b-file row through 27 satisfies this condition.  The stronger
normalization check
\[
 \frac{a(n)-3A033323(n)}{48}\in\mathbb Z_{\geq0}
\]
also passes for each published row $n=23,\ldots,27$, giving respectively
\[
\begin{aligned}
&17838020877573,\quad 79036635370667,\quad 348751024443185,\\
&1534271679749333,\quad 6725942965307646.
\end{aligned}
\]
These are derived orbit counts used as transcription checks, not additional
campaign outputs.  The planar inputs in these five checks are
$A033323(23),\ldots,A033323(27)=$
\[
3150044696,\quad 7994665480,\quad 20209319824,\quad
50942982080,\quad 127962421824,
\]
respectively; they are the Bennett-Wood table values after the conversion
specified above and coincide with the corresponding OEIS A033323 terms~
\cite{BennettWoodEtAl1998,OEISA033323}.

Second, over the newly computed block the planar share
$3A033323(n)/a(n)$ decreases from $0.010275\%$ at $n=19$ to $0.001927\%$ at
$n=22$ (the intermediate values are $0.005879\%$ and $0.003366\%$).  The
successive ratios $a(n)/a(n-1)$ also decrease throughout the currently
completed block, from 4.539875 at $n=19$ to 4.383798 at $n=27$.  These are
descriptive observations on a short exact range.  They neither establish
monotonicity beyond 27 nor estimate a critical exponent or growth constant.
No new conjecture is proposed.

\section{Conclusion}

The literature audit first identified $n=19,20,21,22$ as the only unresolved
indices intervening between the OEIS prefix through $a(18)$ and the normalized
published values at $n=23,\ldots,27$.  The exhaustive orbit-reduced computation
reported here determines those four values.  Exact dimensional reconstruction,
old-prefix replay, small direct and explicit-orbit oracles, contact-class and
axial checks, manifest validation, and a separate direct Java reproduction at
$n=19$ give complementary evidence.  The boundary is
therefore precise: the new campaign ends at 22, the documented sequence ends
at 27, and, within the evidence assembled in the audit of 26 August 2026, 28
remains unresolved.

\section*{Data availability}

The source, Makefile, frozen campaign manifest and result table, derived
sequence table, canonical b-file, validators, independent jOEIS harness,
completion report, and manuscript source are prepared for inclusion in a
synchronized public release.  The repository and archive locators are listed
in the References.  The GitHub URL identifies the intended public release
repository; the Zenodo DOI is a draft placeholder that will be replaced after
the first GitHub review and GitHub--Zenodo synchronization.  Neither location
is claimed to be public at the time of this draft.

\section*{AI use disclosure}

Large language models were used as research, code-review,
validation-planning, artifact-auditing, and drafting assistants during the
A049230 project.  The preserved substantive records identify OpenAI Codex and
review systems from the ChatGPT, Claude, DeepSeek, Gemini, Grok, Kimi, Lumo,
Meta, MiniMax, Qwen, and Z families.  They are named at family level where
exact model-version metadata were not consistently available; no missing
version identifier has been reconstructed from inference.

\textbf{Mathematical analysis.}
LLMs assisted in restating the dimensional decomposition, challenging the
cubic-orbit and planar-stabilizer arguments, locating step/monomer and
first-step normalization risks, and identifying exact checks from contact
polynomials and related sequences.  Accepted statements were checked against
the OEIS snapshots, primary papers, exact arithmetic, the direct rooted oracle,
and the explicit 48-image orbit oracle.  Conversational assertions were not
treated as mathematical evidence.  In the second external manuscript round,
several independent reviews identified the omitted transverse-axis ordering
condition in the written canonical rule.  The correction in Proposition~
\ref{prop:canonical} was accepted only after direct reconciliation with the
frozen C source and the exhaustive small-$n$ orbit oracle.

\textbf{Algorithm and code.}
OpenAI ChatGPT/Codex and Claude-family systems supplied the final corrective
code reviews; systems from the other named families contributed substantive
review artifacts at earlier stages.  Their work included static or dynamic
code review, adversarial test design, persistence analysis, parser and
output-boundary review, and artifact inspection.  The final numerical claims
rest on the frozen C17 program and build recipe, the completed manifest and result table,
known-term replay, cross-thread and split-depth tests, bounded direct and orbit
oracles, axial-data checks, the arithmetic derivator, and the separate jOEIS
reproduction at $n=19$.  Shared-backend checks and unavailable tool evidence
were not promoted to independent proofs.

\textbf{Project workflow.}
AI systems assisted in organizing bounded review rounds, comparing findings,
maintaining evidence boundaries, and preparing operator-facing runbooks.  They
did not launch or operate the official computation.  Carlo Corti authorized,
started, monitored, and completed the campaign and made the scientific and
editorial decisions.

\textbf{Manuscript.}
Codex assisted with literature research, drafting, claim-to-source auditing,
arithmetic consistency checks, and LaTeX diagnostics.  Every numerical result
retained in this revision was checked against a result table, the canonical
b-file, the completed manifest, validator output, exact arithmetic, or a cited
primary source.  The final claims depend on those artifacts and sources, not
on LLM conversation.

Carlo Corti takes full scientific responsibility for the mathematical
statements, algorithm, implementation, numerical results, validation
boundaries, and conclusions of this paper.

\clearpage
\begin{thebibliography}{99}
\sloppy
\raggedright

\bibitem{FisherHiley1961}
M.~E. Fisher and B.~J. Hiley,
\newblock Configuration and free energy of a polymer molecule with solvent
interaction,
\newblock \emph{J. Chem. Phys.} \textbf{34} (1961), 1253--1267.
\url{https://doi.org/10.1063/1.1731729}.

\bibitem{NemirovskyEtAl1992}
A.~M. Nemirovsky, K.~F. Freed, T. Ishinabe, and J.~F. Douglas,
\newblock Marriage of exact enumeration and $1/d$ expansion methods: lattice
model of dilute polymers,
\newblock \emph{J. Stat. Phys.} \textbf{67} (1992), 1083--1108.
\url{https://doi.org/10.1007/BF01049010}.

\bibitem{Crippen2000}
G.~M. Crippen,
\newblock Enumeration of cubic lattice walks by contact class,
\newblock \emph{J. Chem. Phys.} \textbf{112} (2000), 11065--11068.
\url{https://doi.org/10.1063/1.481746}.

\bibitem{BennettWoodEtAl1998}
D. Bennett-Wood, I.~G. Enting, D.~S. Gaunt, A.~J. Guttmann, J.~L. Leask,
A.~L. Owczarek, and S.~G. Whittington,
\newblock Exact enumeration study of free energies of interacting polygons and
walks in two dimensions,
\newblock \emph{J. Phys. A: Math. Gen.} \textbf{31} (1998), 4725--4741.
\url{https://doi.org/10.1088/0305-4470/31/20/010}.

\bibitem{LeeEtAl2012}
J.~H. Lee, S.-Y. Kim, and J. Lee,
\newblock Exact partition function zeros of a polymer on a simple cubic
lattice,
\newblock \emph{Phys. Rev. E} \textbf{86} (2012), 011802.
\url{https://doi.org/10.1103/PhysRevE.86.011802}.

\bibitem{HsiehEtAl2016EPJ}
Y.-H. Hsieh, C.-N. Chen, and C.-K. Hu,
\newblock Exact partition functions of interacting self-avoiding walks on
lattices,
\newblock \emph{EPJ Web Conf.} \textbf{108} (2016), 01005.
\url{https://doi.org/10.1051/epjconf/201610801005}.

\bibitem{HsiehEtAl2016CPC}
Y.-H. Hsieh, C.-N. Chen, and C.-K. Hu,
\newblock Efficient algorithm for computing exact partition functions of
lattice polymer models,
\newblock \emph{Comput. Phys. Commun.} \textbf{209} (2016), 27--33.
\url{https://doi.org/10.1016/j.cpc.2016.08.006}.

\bibitem{HuangEtAl2022}
S.-S. Huang, Y.-H. Hsieh, and C.-N. Chen,
\newblock Exact enumeration approach to estimate the theta temperature of
interacting self-avoiding walks on the simple cubic lattice,
\newblock \emph{Polymers} \textbf{14} (2022), 4536.
\url{https://doi.org/10.3390/polym14214536}.

\bibitem{ShiraiSakumichi2023}
N.~C. Shirai and N. Sakumichi,
\newblock Solvent-induced negative energetic elasticity in a lattice polymer
chain,
\newblock \emph{Phys. Rev. Lett.} \textbf{130} (2023), 148101.
\url{https://doi.org/10.1103/PhysRevLett.130.148101}.

\bibitem{OEISA049230}
N.~J.~A. Sloane, P. Hadjicostas, S.~A. Irvine, et al.,
\newblock The On-Line Encyclopedia of Integer Sequences, A049230,
\newblock \url{https://oeis.org/A049230}, accessed 26 August 2026.

\bibitem{OEISA033323}
N.~J.~A. Sloane, P. Hadjicostas, S.~A. Irvine, C. Corti, et al.,
\newblock The On-Line Encyclopedia of Integer Sequences, A033323,
\newblock \url{https://oeis.org/A033323}, accessed 26 August 2026.

\bibitem{JoeisRepository}
S.~A. Irvine et al.,
\newblock jOEIS: Java programs for integer sequences,
\newblock \url{https://github.com/archmageirvine/joeis}, accessed 26 August
2026.

\bibitem{CortiPackage2026}
C. Corti,
\newblock \emph{A049230: Exact simple-cubic-lattice two-contact enumeration},
\newblock planned GitHub repository and Zenodo archive, 2026.
\url{https://github.com/carcorti/A049230}.
\url{https://doi.org/10.5281/zenodo.xxxxxxxx}.
The Zenodo DOI placeholder will be replaced after the first GitHub review.

\end{thebibliography}

\end{document}
